\documentclass[10pt,a4paper]{amsart}
\usepackage[margin=19mm]{geometry}
\usepackage{amsmath,amssymb,mathtools}
\usepackage[scr=rsfs,cal=euler]{mathalfa}
\usepackage{xfrac}
\usepackage[unicode,hidelinks,bookmarks=false]{hyperref}
\newcommand{\ie}{i.e.\ }
\newcommand{\cf}{cf.\ }
\newcommand{\resp}{respectively\ }
\newcommand{\Subsection}{\subsection}
\providecommand{\nobreakdash}{\nobreak\hbox{-}}
\setlength{\emergencystretch}{2em}
\multlinegap0pt
\usepackage{amssymb}
\usepackage{mathtools}
\usepackage{xcolor}
\usepackage{float}
\usepackage[shortlabels]{enumitem}
\usepackage{xfrac}
\usepackage{tikz}
\newtheorem{theorem}{Theorem}
\newtheorem{proposition}{Proposition}
\newtheorem{remark}{Remark}
\newcommand{\D}{\mathbb{{D}}}
\newcommand{\R}{\mathbb{{R}}}
\title{$A^p_\alpha$ classes in the Dirichlet range: inner-outer factorization, Carleson measures and weak products}
\author{Alberto Dayan, Adri\'an Llinares, Miguel Monsalve-L\'opez}
\date{Source: arXiv:2604.21347v2 (CC BY 4.0)}
\begin{document}
\maketitle

\noindent\emph{Source and licence.} $A^p_\alpha$ classes in the Dirichlet range: inner-outer factorization, Carleson measures and weak products. Version arXiv:2604.21347v2 (CC BY 4.0), distributed by arXiv under the Creative Commons Attribution 4.0 International licence, http://creativecommons.org/licenses/by/4.0/, as recorded in the arXiv licence metadata for this exact version and shown on the abstract page https://arxiv.org/abs/2604.21347v2. Full licence notice: \texttt{licenses/arxiv-2604.21347-CC-BY-4.0.txt}.\par\smallskip

\noindent\textit{Editorial scope.} Counted unit: the article's theorem theo:NonLinear (source0.tex lines 197--200). The article's complete original proof of it is delivered with it (source0.tex lines 461--529, one \texttt{proof} environment of 69 lines, which the article places away from the statement). No mathematics is written, repaired or re-derived: the statement and the proof are the article's own lines, in the article's own notation, and no retained line is substituted or re-flowed.  Retained uncounted as the source-local context the proof needs: lines 168--171; lines 224--227; lines 367--378; lines 442--457.  Nothing was omitted from any retained span, so the package carries no omission record.  No retained line contains a citation, so the wrapper carries no bibliography and no bibliography entry.  No retained line contains a figure environment, an image-inclusion command or drawing code, so no image is delivered and the package declares no asset dependency.  Editorial formatting only, in addition: an AMS \texttt{amsart} preamble that restates the article's own declarations and macros used by the retained lines, namely the environments \texttt{theorem}, \texttt{proposition}, \texttt{remark} on separate counters, together with the standard packages those lines need.  The retained environments print in this document as Theorem 1, Proposition 1, Remark 1 in order of appearance, whereas the article prints the same items with the numbering of its own full document.  The complete native source of the article as distributed in the arXiv e-print for this exact version is bundled unchanged as \texttt{sources/199019/source0.tex}.
            \begin{equation}
                \label{eqn:nonvanishing}
                f\in A^p_\alpha\iff f^\frac{p}{2}\in A^2_\alpha.
            \end{equation}

        \begin{theorem}
        \label{theo:p/2}
            Let $0<\alpha<1$, $p>0$, and $f=\Theta h$ a function in $H^p$, $\Theta$ and $h$ being its inner and outer factors, respectively. Then $f$ belongs to $A^p_\alpha$ if and only if $\Theta h^\frac{p}{2}$ is in $A^2_\alpha$. More generally, $f \in A^p_\alpha$ if and only if $\Theta h^{\frac{p}{q}} \in A^q_\alpha$, where $q$ is another positive exponent.
        \end{theorem}

        \begin{proposition} \label{Prop:BS}
            Let $a \ge 0$, $b \in \R$ and $0<\alpha<1$. Consider the function
                \[
                    f_{a, b}(z) := (1-z)^b\exp\left(-a\frac{1+z}{1-z}\right).
                \]
                Then:
                \begin{enumerate}
                    \item \label{Propi)} $f_{a, b}$ belongs to $A_\alpha^2$ if $b > \frac{1}{2} - \alpha$.
                    
                    \item \label{Propii)} If $b>-a$ and $b\le\frac{1}{2}-\alpha$, then $f$ does not belong to $A^p_\alpha$.
                \end{enumerate}
        \end{proposition}

        \begin{remark} \label{rem:power_est}
            Using a change of variables and the monotone convergence theorem, we can show that
                \begin{align*}
                    \lim_{j \to \infty} j^{1 + \alpha} \int_0^{+\infty} \dfrac{u^\alpha}{(u + 1)^{2(jb + \delta + \alpha) - 1}} e^{-2jau} du & = \lim_{j \to \infty} \dfrac{1}{(2a)^{1 + \alpha}} \int_0^{+\infty} \dfrac{t^{\alpha}}{ \left( \frac{t}{2ja} + 1 \right)^{2(jb+\delta + \alpha)-1}} e^{-t} \, dt \\
                        & = \dfrac{1}{(2a)^{1+\alpha}}\int_0^{+\infty} t^\alpha e^{-t (1 + b/a)} \, dt = \dfrac{\Gamma (\alpha + 1)}{\big(2 (a+b) \big)^{1+\alpha}}
                \end{align*}
                if $a + b > 0$. Thus we have that
                \[
                    \|f_{ja, jb+\delta}\|_{\alpha,2} \simeq j^{\frac{1 - \alpha}{2}} 2^{jb + \delta} \sqrt{B \left( jb + \delta + \alpha - \dfrac{1}{2}, \dfrac{1}{2} \right)} \simeq j^{\frac{1}{4} - \frac{\alpha}{2}} 2^{jb}.
                \]
                Since $\left|\binom{\frac{p}{2}}{j}\right|\simeq \frac{1}{j^{1+\frac{p}{2}}}$ when $j \to \infty$, we will have that
                \[
                    \left|\binom{\frac{p}{2}}{j}\right| \dfrac{\|f_{ja, jb+\delta}\|_{\alpha,2}}{2^{jb}} \simeq j^{-\frac{3}{4} - \frac{\alpha + p}{2}}
                \]
                which is summable when $\alpha + p > \frac{1}{2}$.
        \end{remark}

        \begin{theorem}
        \label{theo:NonLinear}
            Let $0<\alpha<1$ and $p \neq 2$. If $\alpha + p > \frac{1}{2}$, there exist two functions $f$ and $g$ in $A^p_\alpha$ such that $f+g$ is not in $A^p_\alpha$.
        \end{theorem}

        \begin{proof}[Proof of Theorem \ref{theo:NonLinear}]
        	Consider the function $f_{1, b}$ as in Proposition \ref{Prop:BS} and 
                \[
                    g_c(z):=2^{b-c}(1-z)^c,
                \]
                where $0<\varepsilon<1-\alpha$ and
                \[
                    c=\frac{\varepsilon-\alpha}{p}.
                \]
                A standard computation using Theorem \ref{theo:p/2} and the Taylor series of $(1-z)^\frac{cp}{2}$ shows that $(1-z)^c$ belongs to $A^p_\alpha$, since $c>-\frac{\alpha}{p}$. On the other hand, Proposition \ref{Prop:BS} yields that $f_{1, b}$ is in $A^p_\alpha$ when 
                \begin{equation}
                \label{eqn:f_1b}
                    b>\frac{1}{p}(1-2\alpha).
                \end{equation}
            
            Let $h(z) := f_{1, b}(z) + g_c(z)$. If $b$ satisfies \eqref{eqn:f_1b}, then $b - c > \frac{1 - \alpha - \varepsilon}{p} > 0$ and hence 
                \[
                    \left| \dfrac{1-z}{2} \right|^{b-c} \left|\exp\left(-\frac{1+z}{1-z}\right)\right| < 1, \quad z \in \mathbb{D}.
                \]
                Thus, the function $h$ is nowhere vanishing in $\D$. Thanks to \eqref{eqn:nonvanishing}, $h$ lies into $A_\alpha^p$ if and only if $h^{p/2}$ belongs to $A^2_\alpha$. Moreover,
            	\[
            		\begin{aligned}
            		\big(h(z)\big)^{p/2} & = 2^{\frac{p(b-c)}{2}} (1-z)^{\frac{cp}{2}}\Bigg(1+ 2^{-(b-c)}(1-z)^{b-c}\exp\bigg(-\frac{1+z}{1-z}\bigg)\Bigg)^{p/2} \\
            		& = 2^{\frac{p(b-c)}{2}} (1-z)^{\frac{cp}{2}} \sum_{j = 0}^\infty \binom{p/2}{j} \frac{\big(f_{1,b-c}(z)\big)^j}{2^{j(b-c)}} \\
            		& = 2^{\frac{p(b-c)}{2}}  \sum_{j = 0}^\infty \binom{p/2}{j} \frac{f_{j,j(b-c) + \frac{cp}{2}}(z)}{2^{j(b-c)}}.
            		\end{aligned}
            	\]
            
            Since $ A^2_\alpha$ is a Hilbert space,
                \[
                    \left\|\big(h(z)\big)^{p/2}\right\|_{\alpha, 2}\le\sum_{j = 0}^\infty \left|\binom{p/2}{j}\right| \frac{\|f_{j,j(b-c) + \frac{cp}{2}}\|_{\alpha, 2}}{2^{j(b-c)}},
                \]
                thus Proposition \ref{Prop:BS} and Remark \ref{rem:power_est} yield that $h \in A^p_\alpha$ if 
                    \begin{equation}
                    \label{eqn:h}
                        b-c+\frac{cp}{2}>\frac{1}{2}-\alpha.
                    \end{equation}
            
            On the other hand, if 
                \begin{equation}
                \label{eqn:cnot}
                    b-c+\frac{cp}{2}=\frac{1}{2}-\alpha.
                \end{equation}
                then
                \[
                    \left\|\big(h(z)\big)^{p/2}\right\|_{\alpha, 2}\ge\left| \frac{p}{2^{b-c+1}}\|f_{1, b-c+\frac{cp}{2}}\|_{\alpha, 2}- \left \| \sum_{j \ne1} \binom{p/2}{j} \frac{f_{j,j(b-c) + \frac{cp}{2}}}{2^{j(b-c)}} \right \|_{\alpha, 2} \right|=\infty,
                \]
                thus $h$ is not in $A^p_\alpha$ in this case.
        
            The argument now distinguishes two cases:
                \begin{itemize}
                
                    \item Suppose that $p>2$. We will choose $b$ such that $f_{1, b}$ is in $A^p_\alpha$ but $h$ is not. To this end, fix $b$ as in \eqref{eqn:cnot}. We are left to check that $f_{1, b}$ is in $A^p_\alpha$: this is the case because inequality \eqref{eqn:f_1b} is equivalent to 
                     \[
                     \varepsilon\left(\frac{1}{p}-\frac{1}{2}\right)>(1-\alpha)\left(\frac{1}{p}-\frac{1}{2}\right),
                     \]
                    which holds since $\varepsilon<1-\alpha$ and $p>2$.
                     
                    \item Now assume that $p<2$, and fix 
                    \[
                    b=\frac{1}{p}(1-2\alpha),
                    \]
                    so that $f_{1, b}$ is not in $A^p_\alpha$. In order to check that $h$ is in $A^p_\alpha$ we are left to show that \eqref{eqn:h} holds. A computation shows that, with the parameters $b$ and $c$ that we fixed, \eqref{eqn:h} is equivalent to 
                    \[
                    \varepsilon\left(\frac{1}{2}-\frac{1}{p}\right)>\left(\frac{1}{2}-\frac{1}{p}\right)(1-\alpha)
                    \]
                    which holds since $\varepsilon<1-\alpha$ and $p<2$.
                \end{itemize}
        \end{proof}

\end{document}
